% Library CV process_fr — transcription of writing-samples/cvs/Moumen_Nermine_procédés.pdf (CV01, 2026-11-11)
\documentclass[10pt]{article}
\usepackage[letter]{nermine-cv}
\begin{document}
\cvheader{Moumen Nermine}{Étudiante en dernière année du cycle ingénieur en Chimie – Génie des procédés, énergie et environnement}{nermine.moumen@insat.ucar.tn}{+21627040414}{linkedin.com/in/nermine-moumen}

\section{FORMATION}
\entry{Institut National des Sciences Appliquées et de Technologie (INSAT)}{2022–Présent}
\summary{Étudiante en dernière année de Chimie Industrielle (spéc. Procédés et Environnement), avec une formation solide en génie des procédés, thermodynamique, mécanique des fluides et opérations unitaires, complétée par la simulation de procédés industriels, la modélisation de systèmes électrochimiques, l'évaluation technico-économique et l'analyse de cycle de vie. À la recherche d'un PFE.}

\section{EXPÉRIENCE PROFESSIONNELLE}
\entry{SEGULA Technologies - Stage en modélisation électrochimique des batteries}{06/2026–08/2026}
\begin{bul}
\item Modélisation d'une cellule \textbf{Li-ion} (NMC811/Graphite-SiOx) via le Single Particle Model (SPM) implémenté sous MATLAB.
\item Résolution des équations aux dérivées partielles par \textbf{FVM} avec intégration temporelle de Crank-Nicolson.
\item Optimisation du maillage via une \textbf{analyse des données} réduisant l'erreur \textbf{RMSE de 202,13 mV à 40,5 mV.}
\item Validation des résultats par une \textbf{approche itérative} et une \textbf{analyse critique} avec les données expérimentales et COMSOL.
\end{bul}
\entry{Control Union - Stage en contrôle qualité des carburants et lubrifiants}{08/2025–09/2025}
\begin{bul}
\item Caractérisation de lubrifiants selon les normes \textbf{ASTM/ISO} : viscosité cinématique et métaux d'usure et additifs par \textbf{ICP-OES}.
\item Rédaction de diagnostics et recommandations dans des \textbf{rapports d'analyse techniques} conformément à la norme \textbf{ISO 17025}.
\end{bul}
\entry{AGIL - Stage en contrôle qualité des carburants}{07/2025–08/2025}
\begin{bul}
\item Réalisation d'analyses physico-chimiques \textbf{(point d'éclair et distillation atmosphérique)} de carburants diesel et Jet A-1.
\end{bul}
\entry{ETAP - Stage en caractérisation des hydrocarbures et traitement des eaux industrielles}{07/2024–08/2024}
\begin{bul}
\item Formation aux techniques de géochimie organique (\textbf{GC}, GC-MS, pyrolyse Rock-Eval) pour la caractérisation d'\textbf{hydrocarbures}.
\item Analyses physico-chimiques des eaux (pH, salinité, bicarbonates, sulfates, chlorures) par chromatographie ionique et ICP-AES.
\end{bul}

\section{PROJETS ACADÉMIQUES}
\project{Modélisation de Piles à Combustible PEM et Électrolyseurs PEM sous COMSOL et ANSYS Fluent}
\begin{bul}
\item Modélisation de piles à combustible : écoulement diphasique, \textbf{effets thermiques}, engorgement de la couche de diffusion.
\item Modélisation d'un \textbf{électrolyseur} PEM 3D : cinétique électrochimique et identification de la réaction limitante (OER).
\end{bul}
\project{Simulation du procédé de production d'oxyde d'éthylène sous Aspen HYSYS V14}
\begin{bul}
\item Modélisation d'un \textbf{réacteur catalytique gaz-solide tubulaire à lit fixe} à partir de ses paramètres de fonctionnement et de sa géométrie avec la \textbf{cinétique de Langmuir-Hinshelwood} de deux réactions.
\item \textbf{Optimisation} des conditions opératoires par analyse de sensibilité, atteignant 96 \% de récupération et \textbf{97,14 \% de pureté}.
\end{bul}
\project{Analyse techno-environnementale de la production d'éthylène : vapocraquage vs déshydratation du bioéthanol}
\begin{bul}
\item Simulation comparative de deux voies de production sous \textbf{Aspen Plus V14}, avec établissement des \textbf{bilans matière et énergie} et des diagrammes de flux \textbf{(PFD)} du procédé, pour évaluer la faisabilité technique de chaque filière.
\item Analyse de \textbf{Cycle de Vie} cradle-to-gate sous \textbf{OpenLCA} pour comparer les \textbf{impacts environnementaux} des deux voies.
\end{bul}
\project{Prédiction des Équilibres Liquide–Liquide pour Systèmes Ternaires d'Extraction}
\begin{bul}
\item Formulation du système thermodynamique LLE et résolution via Rachford-Rice et \textbf{continuation par longueur d'arc}.
\item Implémentation du workflow en Python, MATLAB et Mathematica (spline/fsolve, SciPy/mpltern, Interpolation/FindRoot).
\end{bul}

\section{COMPÉTENCES TECHNIQUES}
\begin{bul}
\item \textbf{Simulation :} Aspen Plus et HYSYS, COMSOL Multiphysics\sep\textbf{Programmation :} \textit{Python, Mathematica, MATLAB}
\item \textbf{Normes :} ISO 9001, ISO 14001, ISO 45001, ISO 17025, ISO 14040\sep\textbf{Données \& Bureautique :} \textit{Excel (avancé), OpenLCA}
\item \textbf{Techniques de Laboratoire :} Chromatographie (HPLC, GC, CCM), Spectrométrie (MS, ICP-OES), Spectroscopie (RMN, FTIR, UV-Vis), Diffraction (XRD), Microscopie (MEB, MET)
\end{bul}

\section{ACTIVITÉS EXTRASCOLAIRES}
{\textbf{IEEE INSAT} — Chef d'organisation de TRSYP 3.0 ; Responsable des activités sociales – MENASYP ; Présidente de l'EMBS.\\
\textbf{THEATRO INSAT} — Membre de l'équipe de danse.\par}

\section{LANGUES ET CERTIFICATS}
\twocol{Langues}{Anglais (Avancé), Français (Bilingue), Arabe (Natif).}{Certificats}{Aspen HYSYS – MATLAB}
\end{document}
